\documentclass[11pt]{article}
\usepackage{amsmath,amssymb,amsthm}
\usepackage{fullpage}
\usepackage{amsmath,amssymb,amsthm}
\usepackage[T1]{fontenc}
\usepackage{hyperref}

\numberwithin{equation}{section}
\newtheorem{theorem}{Theorem}[section]
\newtheorem{lemma}[theorem]{Lemma}
\newtheorem{remark}[theorem]{Remark}
\newcommand{\R}{\mathbb R}
\newcommand{\E}{\mathbb E}
\newcommand{\PP}{\mathbb P}
\newcommand{\norm}[1]{\left\|#1\right\|}
\newcommand{\diver}{\operatorname{div}}
\newcommand{\KL}{\operatorname{KL}}
\newcommand{\supp}{\operatorname{supp}}
\newcommand{\Sp}{\mathbb S^{d-1}}

\newcommand{\fptitle}{Review of Problem 2 -- ProofCouncil}
\newcommand{\reviewtag}{Reviewer 1}
\usepackage{fp-review}

\begin{document}
\maketitle

\section{Review}

The paper gives the correct solution of the requested minimax rate (even without log factors) using novel ideas for estimation and lower bound alternatives, compared to the published literature. The estimation relies on a plugin approach for the explicit reconstruction,  norm bounds for the transfer operator, which nicely reveal the intrinsic geometry of the problem, and detailed further arguments to derive the correct convergence rate. The error decomposition is different from the human proof and simplifies some of their upper bounds, in particular getting rid of the log factor. The lower bound relies on spherical perturbations in the image domain of the underlying nonlinear forward operator comparable to the human proof solution. 

\paragraph{Correctness.}

The paper seems completely correct. Although some parts are intrinsically technical, the proofs can be followed well and there are no obvious gaps. By defining the right objects and structuring the proofs the comprehension is facilitated. The arguments are mostly novel, sometimes intriguing or even elegant. 

\begin{itemize}
	\item  1: essentially flawless mathematical argument
	\item  0.5: minor errors that are easily corrected
	\item  0: blank solution, or acknowledged partial result
	\item  $-0.5$: major errors
	\item  $-1$: proposed proof makes no progress towards the result.
\end{itemize}
Note that points should be docked for wrong statements/ solutions.

\noindent Score: 1

\paragraph{Exposition.}

The exposition is impressively good. The general discussions in Sections 3.3 and 3.4 are appreciated. A structural improvement for reading would be to start the upper bound with the definition of the estimator and the basic error decomposition, which is now somewhat hidden in Sections 5.1, 5.4 and 5.5. This is, however, no serious drawback.

\begin{itemize}
	\item  1: well-written and clear solution.
	\item  0: argument that is possible to follow with serious effort
	\item  $-1$: very poor exposition
\end{itemize}

\noindent Score: 1

\paragraph{Attributions.}

At a general level, the paper cites well references for the tools used, e.g. from analysis and statistics. A minor issue is that concrete technical statements, which are proved, are not accompanied by a citation of closely related result (e.g. Lemma 4.3 is partly in reference [9], which is only referred to at large way before). The selection of related work is somewhat selective, missing authors like Andrew Stuart, [name redacted] or the numerical analysis community who study the Darcy inverse problem intensively, although not from a pure statistical minimax point of view.


\begin{itemize}
	\item  1: solutions contains fully appropriate attributions
		to related work and theorems being used.
	\item  0: moderate issues with missing attributions or hallucinated references
	\item  $-1$: blatant plagiarism or hallucinations
\end{itemize}

\noindent Score: 

\newpage
\section{Submission with inline comments}


\begin{abstract}
Let $u_f$ solve $\diver(f\nabla u)=1$ in the unit ball $D\subset\R^d$,
$d\ge2$, with zero Dirichlet data, and suppose $u_f$ is observed in
Gaussian white noise of level $n^{-1/2}$. We determine the minimax rate
for recovering the conductivity $f$ in squared $L^2(D)$ loss over a
small H\"older--Zygmund ball $\{\|f-1\|_{C^\alpha(D)}\le r\}$,
$\alpha>1$: it is $n^{-2\alpha/(2\alpha+d+3)}$, with no logarithmic
loss. The upper bound inverts the transport equation
$\nabla u\cdot\nabla f+(\Delta u)f=1$ along the gradient flow of the
smoothed data; a bias cancellation for smoothed products
and a variance gain of one power of the bandwidth from integrating the
differentiated noise along characteristics drive the rate. The lower
bound prescribes exact solutions oscillating in the $d-1$ angular
directions, recovers their conductivities through the same transport
inverse, and applies Fano's inequality.
\end{abstract}

\section{Introduction}

\subsection{The problem}

Let $D=\{x\in\R^d:|x|<1\}$ be the open unit ball, $d\ge2$, and write
$B_\rho=\{x:|x|<\rho\}$. A \emph{conductivity} is a bounded measurable
function $f$ on $D$ bounded below by a positive constant. For such
$f$, Darcy's problem
\begin{equation}\label{eq:pde}
 \diver(f\nabla u)=1\ \text{ in }D,\qquad u=0\ \text{ on }\partial D,
\end{equation}
has a unique weak solution $u_f\in H^1_0(D)$, by the Lax--Milgram
theorem. The conductivity is unknown, and $u_f$ is observed in Gaussian
white noise of level $n^{-1/2}$:
\begin{equation}\label{eq:model}
 Y_n=u_f+n^{-1/2}\mathbb W .
\end{equation}
Concretely, for every $\varphi\in L^2(D)$ one observes
$Y_n(\varphi)=\langle u_f,\varphi\rangle+n^{-1/2}\mathbb W(\varphi)$,
where $\mathbb W$ is the isonormal Gaussian process on $L^2(D)$: the
variables $\mathbb W(\varphi)$, $\varphi\in L^2(D)$, are jointly
centred Gaussian with covariance
$\E\,\mathbb W(\varphi)\mathbb W(\psi)=\langle\varphi,\psi\rangle$.
We write $P_f$ for the law of $Y_n$ and $\E_f$ for the corresponding
expectation. An \emph{estimator} $\widehat f_n$ is any measurable
function of the observed process with values in $L^2(D)$.

\paragraph{H\"older--Zygmund spaces.}
For $s>0$ let $C^s(\R^d)=B^s_{\infty,\infty}(\R^d)$, with norm
$\norm F_{C^s}=\sup_{i\ge0}2^{si}\norm{F_i}_\infty$, where
$F=\sum_{i\ge0}F_i$ is the inhomogeneous Littlewood--Paley
decomposition ($F_i$ frequency-localized to
$\{2^{i-1}\le|\xi|\le2^{i+1}\}$ for $i\ge1$) \cite[Section 4.3]{GN},
\cite{Triebel}. For noninteger $s$ this is the classical H\"older
space; for integer $s$ it is the Zygmund class, strictly larger than
the space of functions with bounded continuous derivatives of order
$s$. The space $C^s(D)$ consists of the restrictions to $D$ of elements
of $C^s(\R^d)$, with the quotient norm
$\norm f_{C^s(D)}=\inf\{\norm F_{C^s(\R^d)}:F|_D=f\}$
\cite[Section 4.3.6]{GN}, normalized so that
$\norm f_{L^\infty(D)}\le\norm f_{C^s(D)}$. Integer indices always
refer to the Zygmund scale, and classical derivative bounds are written
explicitly as $\norm{D^kF}_\infty$. For an integer $k<s$, Bernstein's
inequality $\norm{D^kF_i}_\infty\le C2^{ki}\norm{F_i}_\infty$
\cite[Lemma 4.3.4]{GN} makes $\sum_iD^kF_i$ converge uniformly, so
every $F\in C^s$ has bounded continuous derivatives of order $k$ with
$\norm{D^kF}_\infty\le C\norm F_{C^s}$ (on $D$, apply this to a bounded
extension). Thus, for $\alpha>1$, $f\in C^\alpha(D)$ has a bounded
gradient and $u\in C^{\alpha+1}(D)$ has bounded second derivatives.

\paragraph{The parameter class.}
For $\alpha>1$ and $r\in(0,1)$ let $\mathcal F(\alpha,r)$ be the set
of all $f\in C^\alpha(D)$ with $\norm{f-\mathbf 1}_{C^\alpha(D)}\le r$,
where $\mathbf 1$ is the constant function $1$. Its members satisfy
$1-r\le f\le1+r$, so they are conductivities. The question is the order,
as $n\to\infty$ for small fixed $r$, of the minimax risk
$\mathcal R_n(\alpha,r)$ defined in the following theorem.

\subsection{Main result}

\begin{theorem}\label{thm:main}
For every $d\ge2$ and $\alpha>1$ there is $r_0>0$ such that, for each
fixed $0<r\le r_0$, there are constants $0<c<C<\infty$ with
\[
 c\,n^{-\frac{2\alpha}{2\alpha+d+3}}
 \ \le\
 \mathcal R_n(\alpha,r)
 =\inf_{\widehat f_n}\sup_{f\in\mathcal F(\alpha,r)}
 \E_f\norm{\widehat f_n-f}_{L^2(D)}^2
 \ \le\
 C\,n^{-\frac{2\alpha}{2\alpha+d+3}}
\]
for all sufficiently large $n$.
\end{theorem}

The threshold $r_0$ depends on $\alpha,d$ only; $c,C$ may depend on
$\alpha,d,r$. Smallness of $r$ is used only in the
upper bound, through Lemma \ref{lem:regularity}, the perturbative
description of $u_f$ that makes it uniformly strictly convex (Hessian
bounded below by a positive multiple of the identity); the lower bound
only needs its alternatives to lie in $\mathcal F(\alpha,r)$.

\subsection{The transport equation, the exponent, and the plan}\label{sec:heuristic}

Expanding the divergence in \eqref{eq:pde} shows that $f$ solves a
\emph{first-order linear transport equation} whose coefficients are
determined by the solution:
\begin{equation}\label{eq:transportdef}
 T_vq:=\diver(q\nabla v)=\nabla v\cdot\nabla q+(\Delta v)\,q,
 \qquad T_{u_f}f=1 .
\end{equation}
For $f\equiv1$ the solution is $u_*(x)=(|x|^2-1)/(2d)$, with
$\nabla u_*=x/d$, $D^2u_*=d^{-1}I$ and $\Delta u_*=1$. For $f$ in a
small ball around $1$, Lemma \ref{lem:regularity} shows that $u_f$ is
a small $C^{\alpha+1}$ perturbation of $u_*$, hence uniformly strictly
convex with outward-pointing gradient on $\partial D$ and a single
interior critical point. The characteristics of $T_{u_f}$ are the
gradient flow lines of $u_f$, which all emanate from that critical
point, and along each of them, in the direction of increasing $u_f$,
$T_{u_f}f=1$ is a scalar linear ODE damped by the coefficient
$\Delta u_f\ge1/2$. Integrating it outward from the critical point
gives an explicit bounded inverse $T_v^{-1}$ (Lemma
\ref{lem:transport}), so that $f=T_{u_f}^{-1}1$ \emph{without any
boundary values of $f$}. Treating \eqref{eq:pde} as a hyperbolic
equation for the coefficient goes back to Richter \cite{Richter}; see
\cite{KSV} for gradient-flow characteristics in a statistical setting.

The exponent arises as follows. Smoothing the data at bandwidth $h$
yields a kernel smoother $V_h\approx u_f$ (Section
\ref{sec:estimator}) whose second derivatives carry noise of
pointwise size $n^{-1/2}h^{-d/2-2}$. A crude bound on the
plug-in $T_{V_h}^{-1}1$ gives a squared noise contribution
$n^{-1}h^{-d-4}$, against a bias of order $h^\alpha$ (itself the result
of a cancellation, Section \ref{sec:residual}), and balancing would
give the denominator $2\alpha+d+4$. The variance is in fact at most
$n^{-1}h^{-d-3}$: $T_v^{-1}$ integrates along characteristics, which
averages out one power of $h$ in the flow direction while leaving the
$d-1$ transverse directions untouched (Lemma \ref{lem:variance}), and
$h^{2\alpha}=n^{-1}h^{-d-3}$ gives $2\alpha+d+3$. The lower bound
matches: solutions $u_\theta=u_*+v_\theta$ oscillating at frequency
$h^{-1}$ in the $d-1$ angular directions only, with $L^2$ size
$h^{\alpha+2}$, have conductivities
$f_\theta-1\approx-T_{u_*}^{-1}\Delta v_\theta$ of $L^2$ size
$h^\alpha$; with $h^{-(d-1)}$ independent bumps, the Kullback--Leibler
budget $nh^{2\alpha+4}\lesssim h^{-(d-1)}$ gives the same $h$.

\emph{Plan.} Section \ref{sec:tools} collects four tools: the
perturbative description of $u_f$ (Lemma \ref{lem:regularity}); the
\emph{product-smoothing cancellation} (Lemma \ref{lem:comm}): for a
kernel smoother $S_h$ with many vanishing moments,
$(S_hF)(S_hG)-S_h(FG)$ is of order $h^{\alpha+1}$, one order better
than $S_hF-F=O(h^\alpha)$, also after up to two derivatives; the
transport inverse $T_v^{-1}$ for uniformly convex $v$ (Lemma
\ref{lem:transport}); and the \emph{variance gain}
$\E\|T_v^{-1}\diver(a\nabla e)\|_2^2\lesssim n^{-1}h^{-d-3}$ for
deterministic $v,a$ and smoothed noise $e$ (Lemma \ref{lem:variance}).
Section \ref{sec:upper} proves the upper bound: the plug-in operator
$T_{V_h}$ is random, so no independence can be exploited; instead an
exact identity \eqref{eq:expansion} expresses the error through the
\emph{deterministic} inverse $T_{u_h}^{-1}$, $u_h$ the smoothed true
solution, applied to a linear noise term plus a quadratic remainder
controlled on a high-probability event. Section \ref{sec:lower}
prescribes the solutions $u_\theta$ and recovers
$f_\theta=T_{u_\theta}^{-1}1$, so that the Kullback--Leibler
divergences are exact Gaussian shifts; Fano's inequality on a
Varshamov--Gilbert packing finishes the proof.

\subsection{Related work, notation and conventions}\label{sec:notation}

The white-noise Darcy model \eqref{eq:model} is studied in Nickl, van
de Geer and Wang \cite{NvdGW}, who obtain, for integer
$\alpha>\max\{d/2+2,\,2d-1\}$ and Sobolev-type classes with boundary
restrictions on $f$, the prediction rate
$n^{-(\alpha+1)/(2\alpha+d+2)}$ for $\E\|\widehat u-u_f\|_{L^2}$,
$\widehat u$ their estimator of $u_f$, and a conductivity bound of the form $n^{-(\alpha-1)/(2\alpha+d+2)}$ for the
expected (non-squared) $L^2$ error via a stability estimate
\cite[Theorems 9--11]{NvdGW}; their lower-bound construction already
uses the unit ball, the source $1$, the baseline $f=1$ and $u_*$.
Posterior contraction for Darcy appears in \cite{GiNi} and in the
monograph \cite{Nickl23}; information-theoretic aspects of nonlinear
inverse problems, including Darcy with a constant background and
transport equations, are discussed in \cite{NP}. We are not aware of
an earlier determination of the minimax rate over
$\mathcal F(\alpha,r)$; as the classes and smoothness ranges differ,
the comparison with the rates above is not literal.
The product-smoothing estimate, the variance gain and the
solution-first alternatives are proved here directly.

Constants $C,c,c',\dots$ are positive and finite, may change from line
to line, and may depend on $\alpha$, $d$, $r$ and on the fixed kernels
and bumps chosen below, but never on $n$, $h$ or the true conductivity;
subscripts indicate additional dependence ($C_A$, $C_k$). ``For large
$n$'' means for $n\ge n_0$ with $n_0$ depending on the same quantities;
$A\asymp B$ means $cB\le A\le CB$. $D^kF$ is the $k$-th derivative of
$F$ (a symmetric $k$-tensor) and $\norm{D^kF}_\infty$ the supremum of
its Euclidean norm; $\nu$ is the outward unit normal of a ball; matrix
inequalities are in the quadratic-form sense and $I$ is the identity.
$\norm\cdot_\infty$ and $\norm\cdot_2$ are the supremum and $L^2$
norms over the domain under discussion ($D$ unless another set is
indicated). For $0<h<1/2$ put $\Omega_h=B_{1-h}$. Polar coordinates
are $x=\rho\omega$, $\rho=|x|$, $\omega\in\Sp$, and $\Delta_{\Sp}$ is
the Laplace--Beltrami operator of the unit sphere. Finally,
$L_n=\sqrt{\log n}$.

\section{Analytic tools}\label{sec:tools}

\subsection{Regularity near the constant conductivity}

\begin{lemma}\label{lem:regularity}
There are $r_0>0$ and $C<\infty$, depending only on $\alpha$ and $d$,
such that for $0<r\le r_0$ and every $f\in\mathcal F(\alpha,r)$,
\begin{equation}\label{eq:regularity}
 \norm{u_f-u_*}_{C^{\alpha+1}(D)}\le Cr,\qquad
 \norm{f}_{C^\alpha(D)}+\norm{\nabla u_f}_{C^\alpha(D)}\le C.
\end{equation}
In particular $u_f$ is classically twice continuously differentiable
on $\overline D$ and is a classical solution of \eqref{eq:pde}, and
\begin{equation}\label{eq:classical}
 \sum_{j=0}^2\norm{D^j(u_f-u_*)}_\infty\le Cr,\qquad
 \norm f_\infty+\norm{\nabla f}_\infty
 +\sum_{j=0}^{2}\norm{D^ju_f}_\infty\le C.
\end{equation}
\end{lemma}

\begin{proof}
Let $\Delta_D^{-1}$ denote the solution operator of the Dirichlet
problem for the Laplacian on $D$. We use the bound
\begin{equation}\label{eq:schauder}
 \norm{\Delta_D^{-1}\diver F}_{C^{\alpha+1}(D)}\le C\norm F_{C^\alpha(D)}
\end{equation}
for vector fields $F$. For noninteger $\alpha$, $\diver F$ lies in the
classical H\"older space $C^{\alpha-1}$, and \eqref{eq:schauder} is
global Schauder theory for the Dirichlet problem: existence and
uniqueness in $C^{2,\gamma}(\overline D)$ (classical functions with
$\gamma$-H\"older second derivatives), the global a priori estimate
and higher boundary regularity \cite[Theorems 6.14, 6.6 and 6.19]{GT}.
At integer $\alpha$ the same linear operator is bounded between the
spaces with two neighbouring noninteger indices, and real interpolation
of H\"older--Zygmund spaces on smooth domains
\cite[Sections 1.6, 4.5 and 5.1--5.2]{Triebel} gives
\eqref{eq:schauder} at the endpoint; the image consists of functions
vanishing on $\partial D$ in every case, since the operator is the same.

Put $a=f-1$, so $\norm a_{C^\alpha(D)}\le r$, and $v=u_f-u_*$. Since
$\Delta u_*=1$, the equation \eqref{eq:pde} for $u_f=u_*+v$ reads
$\Delta v=-\diver(a\nabla u_*+a\nabla v)$ with $v=0$ on $\partial D$,
i.e. $v=-\Delta_D^{-1}\diver(a\nabla u_*+a\nabla v)$. Since pointwise
multiplication is bounded $C^\alpha\times C^\alpha\to C^\alpha$ for
$\alpha>0$ (the algebra property),
$\norm{a\nabla v}_{C^\alpha}\le C\norm a_{C^\alpha}\norm{v}_{C^{\alpha+1}}
\le Cr\norm v_{C^{\alpha+1}}$, so by \eqref{eq:schauder} the right-hand
side is a contraction of the space of $C^{\alpha+1}(D)$ functions with
zero boundary trace when $r\le r_0$ is small, and its fixed point
satisfies $\norm v_{C^{\alpha+1}}\le C\norm{a\nabla u_*}_{C^\alpha}\le
Cr$. Since $\alpha+1>2$, the fixed point has bounded continuous second
derivatives, so it is a classical solution and coincides with the
unique weak solution; this proves
\eqref{eq:regularity}, and \eqref{eq:classical} follows from the
embedding recalled in Section 1.1 ($\alpha>1$ and $\alpha+1>2$).
\end{proof}

\subsection{Kernel smoothing and the product-smoothing cancellation}\label{sec:smoothing}

Fix an integer $L>\alpha+3$ and a real kernel $K\in C_c^\infty(B_1)$
with
\begin{equation}\label{eq:kernel}
 \int K=1,\qquad \int x^\beta K(x)\,dx=0\quad(1\le|\beta|<L).
\end{equation}
Such a kernel exists, e.g. a suitable finite linear combination of
dilates of one smooth bump, and need not be nonnegative; for
background on higher-order kernels see \cite[Section 1.2.2]{Tsybakov}.
Set
$K_h(x)=h^{-d}K(x/h)$ and $S_hF=K_h*F$. Because
$\supp K$ is a compact subset of the open ball $B_1$, for $x$ in a
neighbourhood of $\overline{\Omega_h}$ the function $K_h(x-\cdot)$ is
supported inside $D$: on $\Omega_h=B_{1-h}$ the smoother only uses
values in $D$, and it applies to the white noise,
$S_h\mathbb W(x)=\mathbb W(K_h(x-\cdot))$, a centred Gaussian field
with smooth sample paths. Three elementary facts are used repeatedly
on $\Omega_h$: derivatives commute with $S_h$, $D^kS_hF=S_hD^kF$
whenever $F$ has bounded continuous classical derivatives through
order $k$, so $\norm{D^kS_hF}_\infty\le\norm K_{L^1}\norm{D^kF}_\infty$;
placing derivatives on the kernel instead gives
$\norm{D^kS_hF}_\infty\le Ch^{-k}\norm F_\infty$; and by
\eqref{eq:kernel}, $S_h$ reproduces polynomials of degree less than
$L$, in particular $S_hu_*=u_*$.

\begin{lemma}[Product-smoothing cancellation]\label{lem:comm}
For $F,G$ in a bounded subset of $C^\alpha(D)$, $\alpha>1$, and
$0<h<1/2$,
\begin{align}
 \norm{S_hF-F}_{L^\infty(\Omega_h)}&\le Ch^\alpha,
       \label{eq:approx}\\
 \norm{D^k\{(S_hF)(S_hG)-S_h(FG)\}}_{L^\infty(\Omega_h)}
       &\le Ch^{\alpha+1-k},\qquad k=0,1,2 .
       \label{eq:comm}
\end{align}
These statements include integer H\"older--Zygmund indices.
\end{lemma}

The point of \eqref{eq:comm} is the exponent $\alpha+1$: each factor is
approximated only to order $h^\alpha$, but
$Q_h(F,G):=(S_hF)(S_hG)-S_h(FG)$ vanishes identically when either
argument is constant, so only increments of both $F$ and $G$ enter.

\begin{proof}
Replace $F,G$ by extensions to $\R^d$ with bounded $C^\alpha$ norms;
the convolutions on $\Omega_h$ are unchanged. Use the inhomogeneous
dyadic decompositions $F=\sum_{i\ge0}F_i$, $G=\sum_{j\ge0}G_j$ of
Section 1.1, so that
\[
 \norm{F_i}_\infty,\norm{G_i}_\infty\le C2^{-\alpha i},
 \qquad \norm{D^\ell F_i}_\infty
       \le C_\ell2^{\ell i}\norm{F_i}_\infty,
\]
and likewise for $G_i$ (dyadic characterization and Bernstein
inequality, \cite[Section 4.3]{GN}, \cite{Triebel}). Let
$2^J\le h^{-1}<2^{J+1}$; blocks with index at most $J$ are ``low
frequency''. Since $Q_h$ is bilinear, we estimate $Q_h(F_i,G_j)$ for
low--low pairs individually and treat the remaining pairs in bulk.

\emph{Low--low pairs.} As $Q_h(F,G)=Q_h(G,F)$, it suffices to take
$i\le j\le J$, and we claim
\begin{equation}\label{eq:lowlow}
 \norm{D^k Q_h(F_i,G_j)}_\infty
 \le C h^L2^i2^{(L-1+k)j}
              \norm{F_i}_\infty\norm{G_j}_\infty .
\end{equation}
The factor $2^i$, rather than $2^{Li}$ or $2^{(L-1)i}2^j$, is the
cancellation. It is seen from the exact identity (expand the product
and use $\int K=1$)
\[
 Q_h(F,G)(x)=-\frac12\iint K(y)K(z)
 [F(x-hy)-F(x-hz)][G(x-hy)-G(x-hz)]\,dy\,dz.
\]
Fix $x$ and regard $q(t)=Q_t(F_i,G_j)(x)$ as a function of the
bandwidth $t\in[0,h]$. By \eqref{eq:kernel} and Taylor expansion,
$S_tP(x)=P(x)+O(t^L)$ for every smooth $P$, so $q^{(l)}(0)=0$ for
$l<L$ and $|q(h)|\le h^L\sup_{t\le h}|q^{(L)}(t)|/L!$. \review{You use $|S_tP(x)-P(x)|\le t^L\|P^{(L)}\|_\infty/L!$ exactly, which is more than needed and would require a proof if exact constants were required.} Differentiate
the identity $L$ times in $t$ by the Leibniz rule. A term in which
both brackets are differentiated, $a\ge1$ times the first and $L-a\ge1$
times the second, is bounded by
$C2^{ai}2^{(L-a)j}\norm{F_i}_\infty\norm{G_j}_\infty
 \le C2^i2^{(L-1)j}\norm{F_i}_\infty\norm{G_j}_\infty$
since $i\le j$. If the first bracket is undifferentiated, bound it by
its first difference, $|F_i(x-ty)-F_i(x-tz)|\le Ct2^i\norm{F_i}_\infty$,
and the second by $C2^{Lj}\norm{G_j}_\infty$; since $t2^j\le1$ the same
bound results, and symmetrically if the second bracket is
undifferentiated (then $2^{Li}\cdot t2^j\le2^i2^{(L-1)j}$). Spatial
derivatives $D^k$ fall on $F_i$ or $G_j$ and cost at most $2^{kj}$ by
Bernstein's inequality, which proves \eqref{eq:lowlow}. Summing over
$i\le j\le J$ and using $\alpha>1$ in the inner sum,
\[
 C h^L\sum_{j\le J}2^{(L-1+k-\alpha)j}
       \sum_{i\le j}2^{(1-\alpha)i}
 \le Ch^L2^{(L-1+k-\alpha)J}\le Ch^{\alpha+1-k}.
\]

\emph{Pairs involving a high frequency.} Put
$F_{\rm lo}=\sum_{i\le J}F_i$, $F_{\rm hi}=F-F_{\rm lo}$, and similarly
$G_{\rm lo},G_{\rm hi}$. Summing the block bounds,
\[
 \begin{aligned}
 \norm{F_{\rm hi}}_\infty&\le Ch^\alpha,&
 \norm{D^\ell F_{\rm lo}}_\infty&\le Ch^{1-\ell}\quad(\ell\ge1),\\
 \norm{D^\ell(S_h-I)F_{\rm lo}}_\infty&\le Ch^{\alpha-\ell}
                  &&(0\le\ell\le2),
 \end{aligned}
\]
and the same for $G$; the last line uses $L>\alpha+3$ and the block
bound $\norm{D^\ell(S_h-I)F_i}_\infty\le C2^{(\ell-\alpha)i}(h2^i)^L$.
For the mixed term write
\begin{align*}
 Q_h(F_{\rm lo},G_{\rm hi})
 &=((S_h-I)F_{\rm lo})\,S_hG_{\rm hi}-[S_h,F_{\rm lo}]G_{\rm hi},\\
 [S_h,F_{\rm lo}]G(x)&:=S_h(F_{\rm lo}G)(x)-F_{\rm lo}(x)\,S_hG(x)
 =\int K_h(x-y)\{F_{\rm lo}(y)-F_{\rm lo}(x)\}G(y)\,dy .
\end{align*}
The first product has $k$-th derivative $O(h^{2\alpha-k})$ for
$k\le2$, by the displayed bounds and
$\norm{D^bS_hG_{\rm hi}}_\infty\le Ch^{-b}\norm{G_{\rm hi}}_\infty$.
Differentiating the commutator $k$ times in $x$ and using
$|F_{\rm lo}(y)-F_{\rm lo}(x)|\le C|y-x|$ together with
$\norm{D^2F_{\rm lo}}_\infty\le Ch^{-1}$ bounds it by
$Ch^{1-k}\norm{G_{\rm hi}}_\infty\le Ch^{\alpha+1-k}$. The term
$Q_h(F_{\rm hi},G_{\rm lo})$ is symmetric, and the high--high term is
$O(h^{2\alpha-k})$ by differentiating the kernels. Since $2\alpha>\alpha+1$, this proves
\eqref{eq:comm}. Finally, \eqref{eq:approx} follows by summing
$\norm{(S_h-I)F_i}_\infty\le C\min\{1,(h2^i)^L\}2^{-\alpha i}$ over $i$.
\end{proof}

\subsection{Inverting the transport equation}\label{sec:transport}

\begin{lemma}[Transport inverse]\label{lem:transport}
Let $\Omega$ be a ball with radius in $[1/2,1]$, and let $v$ be
smooth on a neighbourhood of $\overline\Omega$ with
\begin{equation}\label{eq:geometry}
 mI\le D^2v\le MI\ \text{ on }\overline\Omega,\qquad
 \nabla v\cdot\nu>0\ \text{ on }\partial\Omega,
 \qquad m=\frac1{2d},\quad M=\frac3{2d}.
\end{equation}
Then $v$ has a unique critical point $x_v$ in $\Omega$, its minimum.
Write $c_v=\Delta v$, so $c_v\ge dm=c_0:=1/2$. Let $\Phi_t$ be the
backward gradient flow of $v$, $\dot\Phi_t=-\nabla v(\Phi_t)$,
$\Phi_0(x)=x$. For smooth $q$ on $\overline\Omega$, the unique bounded
classical solution of $T_vp=q$ on $\Omega$ is
\begin{equation}\label{eq:flowinverse}
 p(x)=T_v^{-1}q(x)=\int_0^\infty J_t(x)\,q(\Phi_t(x))\,dt,\qquad
 J_t(x)=\exp\Big\{-\int_0^t c_v(\Phi_s(x))\,ds\Big\}.
\end{equation}
This solution is smooth, $T_v^{-1}$ extends to a bounded operator on
$L^2(\Omega)$, and
\begin{align}
 \norm{T_v^{-1}q}_\infty&\le c_0^{-1}\norm q_\infty,
 &\norm{T_v^{-1}q}_2&\le 2c_0^{-1}\norm q_2,\label{eq:inversebounds}\\
 \norm{\nabla p}_\infty
 &\le (c_0+m)^{-1}
       \big(\norm{\nabla q}_\infty+
          \norm{\nabla c_v}_\infty\norm p_\infty\big).
          \label{eq:gradientinverse}
\end{align}
\end{lemma}

\begin{proof}
Strict convexity and the outward boundary drift give a unique interior
minimum $x_v$ of $v$ on $\overline\Omega$; the backward flow points
into $\Omega$ at the boundary, stays in $\Omega$, and converges to
$x_v$. Along a trajectory, $T_vp=q$ reads
$\frac{d}{dt}p(\Phi_t)=-\nabla v\cdot\nabla p(\Phi_t)
=c_vp(\Phi_t)-q(\Phi_t)$, i.e.
$\frac{d}{dt}\{J_tp(\Phi_t)\}=-J_tq(\Phi_t)$. Integrating over
$t\in[0,\infty)$ gives \eqref{eq:flowinverse} for every bounded
solution, because $J_tp(\Phi_t)\to0$ as $J_t\le e^{-c_0t}$; this is
uniqueness, and $\int_0^\infty e^{-c_0t}\,dt=c_0^{-1}$ gives the
supremum bound. Conversely, differentiating under the integral sign in
\eqref{eq:flowinverse} shows that it defines a smooth solution with
bounded derivatives: this is legitimate because the operator norm of
$D\Phi_t$ is at most $e^{-mt}$ by \eqref{eq:geometry}, higher
derivatives of $\Phi_t$ are bounded by
polynomials in $t$, and every derivative of $J_t$ is bounded by a
polynomial in $t$ times $J_t\le e^{-c_0t}$. For the $L^2$ bound,
Liouville's formula \review{A reference or an intermediate step would be appropriate.} gives
$\det D\Phi_t=\exp\{-\int_0^t\Delta v(\Phi_s)\,ds\}=J_t$, so the change
of variables $y=\Phi_t(x)$ yields
\[
 \int_\Omega J_t(x)^2|q(\Phi_t(x))|^2\,dx
 =\int_{\Phi_t(\Omega)}J_t(\Phi_t^{-1}(y))|q(y)|^2\,dy
 \le e^{-c_0t}\norm q_2^2,
\]
and Minkowski's integral inequality gives
$\norm p_2\le\int_0^\infty e^{-c_0t/2}\,dt\,\norm q_2=2c_0^{-1}\norm q_2$.
By density, $T_v^{-1}$ extends to $L^2(\Omega)$ with the same bound.

For the gradient, differentiate $T_vp=q$:
$\nabla v\cdot\nabla(\nabla p)+(c_vI+D^2v)\nabla p=\nabla q-p\nabla c_v$.
This is a transport equation for the vector $w=\nabla p$ along the same
flow, with the symmetric zero-order coefficient
$\Lambda=c_vI+D^2v\ge(c_0+m)I$ and forcing $\Psi=\nabla q-p\nabla c_v$.
With $w(t)=\nabla p(\Phi_t(x))$ we have $\dot w=\Lambda w-\Psi$, hence
$\frac d{dt}|w|^2=2w\cdot\Lambda w-2w\cdot\Psi\ge2(c_0+m)|w|^2-2|w||\Psi|$,
so that $\frac d{dt}\{e^{-(c_0+m)t}|w(t)|\}\ge-e^{-(c_0+m)t}|\Psi(t)|$
wherever $w\ne0$. Integrating from $0$ to $T$ and letting $T\to\infty$
(the term $e^{-(c_0+m)T}|w(T)|$ vanishes because $w$ is bounded) gives
$|\nabla p(x)|\le\int_0^\infty e^{-(c_0+m)t}|\Psi(\Phi_t(x))|\,dt$,
which is \eqref{eq:gradientinverse}.
\end{proof}

\subsection{A variance gain from integration along characteristics}

\review{$e$ is used throughout and the definition should be in a numbered display, which can later be referred to.}
The forcing $\diver(a\nabla e)$ involves second derivatives of the
smoothed noise $e=n^{-1/2}S_h\mathbb W$, of pointwise variance
$n^{-1}h^{-d-4}$. Integration along characteristics gains a factor
$h$: the forcing is correlated over distance $2h$ while the flow moves
at speed $|\nabla v|\asymp|x-x_v|$, so a trajectory spends time of
order $h/|x-x_v|$ in a correlation ball, and
$\int_{\Omega_h}\min\{1,h/|x-x_v|\}\,dx\asymp h$ for $d\ge2$: the
critical point costs nothing in $L^2$.

\begin{lemma}[Variance gain]\label{lem:variance}
Let $0<h<1/2$, and let $v$ be deterministic and satisfy the
hypotheses of Lemma \ref{lem:transport} on $\Omega=\Omega_h$.
Let $a$ be deterministic and classically continuously differentiable
on $\overline{\Omega_h}$, with
$\norm a_{L^\infty(\Omega_h)}+\norm{\nabla a}_{L^\infty(\Omega_h)}\le A_0$.
For $e=n^{-1/2}S_h\mathbb W$,
\begin{equation}\label{eq:sharpvariance}
 \E\norm{T_v^{-1}\diver(a\nabla e)}_{L^2(\Omega_h)}^2
 \le C A_0^2 n^{-1}h^{-d-3},
\end{equation}
where $C$ depends only on $d$ and the fixed kernel $K$.
\end{lemma}

\begin{proof}
Put $Z=\diver(a\nabla e)=a\Delta e+\nabla a\cdot\nabla e$, a continuous
Gaussian field on $\overline{\Omega_h}$; the flow formula
\eqref{eq:flowinverse} applies to it by uniform approximation with
smooth forcings. Since $D^ke(z)=n^{-1/2}\mathbb W(D^kK_h(z-\cdot))$,
the white-noise kernel of $Z(z)$ is
$n^{-1/2}\{a(z)\Delta K_h(z-\cdot)+\nabla a(z)\cdot\nabla K_h(z-\cdot)\}$,
a function supported in the ball of radius $h$ around $z$ with supremum
at most $CA_0n^{-1/2}h^{-d-2}$ (this uses the classical bound on $a$
and $\nabla a$). Hence, by scaling,
\begin{equation}\label{eq:covariance}
 |\E Z(z)Z(z')|
 \le C A_0^2 n^{-1}h^{-d-4}{\bf1}_{\{|z-z'|\le2h\}} .
\end{equation}
Fix $x\in\Omega_h$ and let $R=|x-x_v|$. Since
$\nabla v(y)=\int_0^1D^2v(x_v+\sigma(y-x_v))\,d\sigma\,(y-x_v)$, the
Hessian bounds \eqref{eq:geometry} give
$e^{-Mt}R\le|\Phi_t(x)-x_v|\le e^{-mt}R$. For $s\ge0$, applying the
upper bound to the trajectory started at $\Phi_t(x)$ and the triangle
inequality,
\[
 |\Phi_{t+s}(x)-\Phi_t(x)|
 \ge (1-e^{-ms})|\Phi_t(x)-x_v|
 \ge (1-e^{-ms})e^{-Mt}R.
\]
The set of $s$ for which the right-hand side is at most $2h$ has
measure $O(he^{Mt}/R)$ within $[0,1]$ and is empty in $[1,\infty)$
unless $he^{Mt}/R\ge c$; consequently
\[
 \int_0^\infty e^{-c_0s}
 {\bf1}_{\{|\Phi_{t+s}(x)-\Phi_t(x)|\le2h\}}\,ds
 \le C\min\{1,he^{Mt}/R\}.
\]
Expanding the square of \eqref{eq:flowinverse}, using
\eqref{eq:covariance} and $J_t\le e^{-c_0t}$, and writing the double
time integral over $\{t'\ge t\}$ with $t'=t+s$ (the other half is
symmetric),
\[
 \E|T_v^{-1}Z(x)|^2
 \le CA_0^2n^{-1}h^{-d-4}
 \int_0^\infty e^{-2c_0t}\min\{1,he^{Mt}/R\}\,dt
 \le CA_0^2n^{-1}h^{-d-4}\min\{1,h/R\};
\]
the last integral is evaluated by splitting at $he^{Mt}=R$ and uses
$2c_0=1>M=3/(2d)$, valid for $d\ge2$. (At $R=0$ use the bound by a
constant.) Finally, uniformly in $x_v\in\Omega_h$,
$\int_{\Omega_h}\min\{1,h/|x-x_v|\}\,dx
 \le Ch^d+Ch\int_h^2\rho^{d-2}\,d\rho\le Ch$ for $d\ge2$, and
integrating the pointwise bound over $x$ proves the claim.
\end{proof}

\section{Upper bound}\label{sec:upper}

\subsection{The estimator}\label{sec:estimator}

Take
\begin{equation}\label{eq:bandwidth}
 h=n^{-1/(2\alpha+d+3)},\qquad V_h=S_hY_n\quad\text{on }\Omega_h,
\end{equation}
and consider $n$ large enough that $h<1/2$. Here
$V_h(x)=Y_n(K_h(x-\cdot))$ is the kernel smoother of the data, a smooth
random function on a neighbourhood of $\overline{\Omega_h}$.

\emph{Gate and plug-in.} If $V_h$ satisfies the Hessian bounds
$mI\le D^2V_h\le MI$ of \eqref{eq:geometry} on $\overline{\Omega_h}$
and $\nabla V_h\cdot\nu\ge1/(4d)$ on $\partial\Omega_h$, set
$\widetilde f_h=T_{V_h}^{-1}1$ on $\Omega_h$, given by Lemma
\ref{lem:transport} with $v=V_h$ and $q\equiv1$; otherwise set
$\widetilde f_h=1$.

\emph{Boundary strip.} We fill the strip $D\setminus\Omega_h$ by radial
polynomial extrapolation, the reflection extension of Hestenes and
Seeley \cite{Hestenes,Seeley} adapted to a shrinking ball. Choose an
integer $b>\alpha$ and constants $a_1,\ldots,a_{b+1}$ with
$\sum_{j=1}^{b+1}a_j(-j)^k=1$ for $0\le k\le b$ (a Vandermonde system).
Define $E_hq=q$ on $\Omega_h$, and for $0<t\le h$, $\omega\in\Sp$,
\begin{equation}\label{eq:extension}
 E_hq((1-h+t)\omega)
       =\sum_{j=1}^{b+1}a_jq((1-h-jt)\omega).
\end{equation}
For small $h$ all radii on the right exceed $1/2$, so only values on
$\Omega_h$ are used, and the moment identities say that $E_h$
reproduces, along each ray, polynomials of degree at most $b$ in the
radial variable. A radial change of variables gives the first bound in
\begin{equation}\label{eq:extensionbounds}
 \norm{E_hq}_{L^2(D)}\le C\norm q_{L^2(\Omega_h)},\qquad
 \norm{E_h(f|_{\Omega_h})-f}_{L^\infty(D)}\le Ch^\alpha
       \norm f_{C^\alpha(D)} .
\end{equation}
For the second bound (at every index, integer or not), apply the
extrapolation error to each dyadic block $F_\ell$ of a bounded
extension $F$ of $f$: Taylor's formula of order $b$ and the moment
identities bound $|E_hF_\ell-F_\ell|$ at distance $t$ from
$\partial\Omega_h$ by $C\min\{1,(t2^\ell)^{b+1}\}2^{-\alpha\ell}$, and
summing over $\ell$ gives $Ct^\alpha$ since $b+1>\alpha$.

\emph{Final estimator.} Our estimator is $E_h\widetilde f_h$, clipped
pointwise to $[1-r,1+r]$; all operations involved are measurable.
Every $f\in\mathcal F(\alpha,r)$ takes values in this interval and
pointwise projection onto it is $1$-Lipschitz, so clipping cannot
increase the error, and the clipped estimator has loss bounded by a
constant on any event.

\subsection{The deterministic residual: bias cancellation}\label{sec:residual}

From now on $f\in\mathcal F(\alpha,r)$ is the true conductivity and
$r\le r_0$ with $r_0$ as in Lemma \ref{lem:regularity}, possibly
decreased below. On $\Omega_h$ put $u_h=S_hu_f$, $f_h=S_hf$ and
$e=V_h-u_h=n^{-1/2}S_h\mathbb W$: $u_h,f_h$ are deterministic
smoothed versions of the truth, and $e$ is the smoothed noise, whose
law does not depend on $f$. The deterministic $u_h$ satisfies
\eqref{eq:geometry} with room to spare when $r_0$ is small: $S_hu_*=u_*$ by the kernel moments, so
$u_h-u_*=S_h(u_f-u_*)$, and by \eqref{eq:classical}
$\sum_{j\le2}\norm{D^j(u_h-u_*)}_\infty\le Cr$; since $D^2u_*=d^{-1}I$
is the midpoint of $[m,M]$ and $\nabla u_*\cdot\nu=(1-h)/d\ge1/(2d)$ on
$\partial\Omega_h$, we may fix $r_0$ so that
\begin{equation}\label{eq:room}
 \tfrac34m\,I+\tfrac14M\,I\le D^2u_h\le\tfrac14m\,I+\tfrac34M\,I,\qquad
 \nabla u_h\cdot\nu\ge\tfrac3{8d}\ \text{ on }\partial\Omega_h .
\end{equation}
Moreover,
\begin{equation}\label{eq:crudederivatives}
 \sum_{j=0}^1\norm{D^jf_h}_\infty
       +\sum_{j=0}^2\norm{D^ju_h}_\infty\le C,
 \qquad
 \norm{D^2f_h}_\infty+\norm{D^3u_h}_\infty\le Ch^{-1}.
\end{equation}
Indeed, by \eqref{eq:classical}, $f$ has a bounded gradient and $u_f$
bounded second derivatives; convolution preserves these bounds, and
placing one further derivative on $K_h$ gives the second bound. In
particular $\norm{f_h}_\infty+\norm{\nabla f_h}_\infty\le C$, which is
the hypothesis on $a$ in Lemma \ref{lem:variance}.

The key deterministic fact is that $f_h$ almost solves the transport
equation of $u_h$. On $\Omega_h$, derivatives commute with $S_h$, so
\[
 f_h\nabla u_h=(S_hf)(S_h\nabla u_f),\qquad
 \diver S_h(f\nabla u_f)=S_h\diver(f\nabla u_f)=S_h1=1
\]
by \eqref{eq:pde}. Hence the residual is a divergence of a
product-smoothing error, and Lemma \ref{lem:comm} with $F=f$,
$G=\partial_iu_f$ (both bounded in $C^\alpha(D)$ by
\eqref{eq:regularity}) and $k=1,2$ gives
\begin{align}
 B_h:=T_{u_h}f_h-1
 &=\diver\{(S_hf)(S_h\nabla u_f)-S_h(f\nabla u_f)\},\label{eq:residual}\\
 \norm{B_h}_\infty+h\norm{\nabla B_h}_\infty&\le Ch^\alpha.\label{eq:residualbound}
\end{align}
There is no boundary term because the convolution kernels used on
$\Omega_h$ are supported inside $D$. This cancellation is why the bias
is $h^\alpha$ rather than $h^{\alpha-1}$, the generic size of
$D^2(u_h-u_f)$.

\subsection{The noise event}\label{sec:noise}

For any fixed $A>0$ there is an event $\mathcal E_n$ of probability at
least $1-C_An^{-A}$ on which, for $k=0,1,2,3$,
\begin{equation}\label{eq:noiseevent}
 \norm{D^ke}_{L^\infty(\Omega_h)}
 \le C_A L_n n^{-1/2}h^{-d/2-k}
 =C_A L_n h^{\alpha+3/2-k},
\end{equation}
the equality being the choice \eqref{eq:bandwidth}. This is a standard
Gaussian supremum bound: $D^ke(x)=n^{-1/2}\mathbb W(D^kK_h(x-\cdot))$
has pointwise standard deviation at most
$\sigma_k=Cn^{-1/2}h^{-d/2-k}$, and the canonical (covariance) distance
between $x$ and $y$ is at most $\sigma_kC\min\{1,|x-y|/h\}$, since the
two kernels differ by $O(|x-y|h^{-d-k-1})$ on a set of measure
$O(h^d)$. Dudley's entropy bound and the Gaussian concentration
inequality for suprema \cite[Section 2.3 and Theorem 2.5.8]{GN} give
\eqref{eq:noiseevent}, the factor $L_n$ absorbing both
$\sqrt{\log(1/h)}$ and the $\sqrt{A\log n}$ deviation.

On $\mathcal E_n$, for large $n$, the gate passes: by \eqref{eq:room}
and $\norm{D^2e}_\infty+\norm{\nabla e}_\infty\le CL_nh^{\alpha-1/2}\to0$,
the random $V_h=u_h+e$ satisfies \eqref{eq:geometry} and
$\nabla V_h\cdot\nu\ge1/(4d)$. Moreover, by \eqref{eq:crudederivatives}
and $\norm{D^3e}_\infty\le CL_nh^{\alpha-3/2}$,
\begin{equation}\label{eq:D3V}
 \norm{D^3V_h}_\infty\le Ch^{-1}\quad\text{on }\mathcal E_n,
 \qquad\text{because }L_nh^{\alpha-1/2}\to0 .
\end{equation}

\subsection{Control of the plug-in error}\label{sec:plugin}

Work on $\mathcal E_n$ with $n$ large. Set $g=\widetilde f_h-f_h$ on
$\Omega_h$, the error of the plug-in relative to the smoothed truth.
Since $T_v$ is linear in $v$, $T_{V_h}=T_{u_h}+T_e$, and by
\eqref{eq:residual},
$T_{V_h}f_h=1+B_h+\diver(f_h\nabla e)$; with $T_{V_h}\widetilde f_h=1$,
the exact equation for $g$ is
\begin{equation}\label{eq:gequation}
 T_{V_h}g=-B_h-\diver(f_h\nabla e),
\end{equation}
and since $g$ is bounded, $g=T_{V_h}^{-1}$ of the right-hand side by
uniqueness in Lemma \ref{lem:transport} (the same remark applies to
\eqref{eq:expansion} below). By \eqref{eq:residualbound} and
\eqref{eq:crudederivatives}, the right-hand side is bounded by
$C(h^\alpha+\norm{D^2e}_\infty+\norm{\nabla e}_\infty)$ and its
gradient by
$C\{h^{\alpha-1}+\norm{D^3e}_\infty+\norm{D^2e}_\infty
+h^{-1}\norm{\nabla e}_\infty\}$, the last term coming from $D^2f_h$.
Applying the supremum bound \eqref{eq:inversebounds} for
$T_{V_h}^{-1}$, then \eqref{eq:gradientinverse} with $v=V_h$, where
$\norm{\nabla c_{V_h}}_\infty\le\sqrt d\,\norm{D^3V_h}_\infty\le Ch^{-1}$
by \eqref{eq:D3V}, and inserting \eqref{eq:noiseevent},
\begin{equation}\label{eq:gbounds}
 \norm g_\infty\le CL_nh^{\alpha-1/2},
 \qquad \norm{\nabla g}_\infty\le CL_nh^{\alpha-3/2}
 \qquad\text{on }\mathcal E_n .
\end{equation}
No independence between $g$ and $e$ is used.

\subsection{Expansion in the deterministic operator and conclusion}

Rewrite \eqref{eq:gequation} with the deterministic transport operator
on the left: $T_{u_h}g=T_{V_h}g-T_eg=-B_h-\diver(f_h\nabla e)
-\diver(g\nabla e)$, so that on $\Omega_h$
\begin{equation}\label{eq:expansion}
 g=-T_{u_h}^{-1}B_h
   -T_{u_h}^{-1}\diver(f_h\nabla e)
   -T_{u_h}^{-1}\diver(g\nabla e).
\end{equation}
The three terms are bias, linear noise and nonlinear remainder; the
inverse is now deterministic, so Lemma \ref{lem:variance} applies to
the middle term. By
\eqref{eq:inversebounds} and \eqref{eq:residualbound}, the first term
has $L^2(\Omega_h)$ norm at most $C\norm{B_h}_2\le Ch^\alpha$. Lemma
\ref{lem:variance}, applied with the deterministic $v=u_h$ (which
satisfies \eqref{eq:geometry} by \eqref{eq:room}) and $a=f_h$ (with
$A_0\le C$ by \eqref{eq:crudederivatives}), bounds the unconditional
second moment of the second term by $Cn^{-1}h^{-d-3}$. On
$\mathcal E_n$, by \eqref{eq:inversebounds}, \eqref{eq:gbounds} and
\eqref{eq:noiseevent}, the last term has $L^2(\Omega_h)$ norm at most
\begin{equation}\label{eq:remainder}
 C\norm{\diver(g\nabla e)}_\infty\le
 C\big(\norm g_\infty\norm{D^2e}_\infty
       +\norm{\nabla g}_\infty\norm{\nabla e}_\infty\big)
 \le CL_n^2h^{2\alpha-1}.
\end{equation}
Finally, $\norm{E_h\widetilde f_h-f}_{L^2(D)}\le
C\norm{\widetilde f_h-f}_{L^2(\Omega_h)}+Ch^\alpha$ by
\eqref{eq:extensionbounds}, and
$\norm{\widetilde f_h-f}_{L^2(\Omega_h)}\le\norm g_2+\norm{f_h-f}_{L^2(\Omega_h)}
\le\norm g_2+Ch^\alpha$ by \eqref{eq:approx}. Using
\eqref{eq:expansion} on $\mathcal E_n$, the bounded loss of the
clipped estimator on the complement, and clipping, the squared risk is
bounded, uniformly in $f\in\mathcal F(\alpha,r)$, by
\begin{equation}\label{eq:riskupper}
 C\{h^{2\alpha}+n^{-1}h^{-d-3}
           +(\log n)^2h^{4\alpha-2}+n^{-A}\}.
\end{equation}
With \eqref{eq:bandwidth}, the first two terms coincide and equal
$n^{-2\alpha/(2\alpha+d+3)}$; since $\alpha>1$, the third is
$(\log n)^2h^{2\alpha-2}\cdot h^{2\alpha}=o(h^{2\alpha})$, so the
nonlinear remainder is negligible; and we may take $A=2$. This proves
the upper bound in Theorem \ref{thm:main}.

\section{Lower bound}\label{sec:lower}

\subsection{Construction of the alternatives}

We need conductivities in $\mathcal F(\alpha,r)$ that are well
separated in $L^2(D)$ while their laws are close. Since
$f\mapsto u_f$ is nonlinear, we prescribe exact solutions $u_\theta$,
so that the Kullback--Leibler divergences are exact Gaussian shifts,
and recover $f_\theta=T_{u_\theta}^{-1}1$ by Lemma
\ref{lem:transport}; the only nonlinear quantity to control is then
the deviation of $f_\theta-1$ from its linearization
$-T_{u_*}^{-1}\Delta(u_\theta-u_*)$.

In a fixed coordinate patch of the sphere $\Sp$ (of dimension
$d-1\ge1$), place $N_h\asymp h^{-(d-1)}$ smooth bumps
$\psi_{1,h},\ldots,\psi_{N_h,h}$ with pairwise disjoint supports of
angular diameter proportional to $h$; they are translates and
rescalings, in the chart, of one fixed nonzero bump. Uniformly in $j$,
\begin{equation}\label{eq:angularbumps}
 \norm{\psi_{j,h}}_{L^2(\Sp)}^2\asymp h^{d-1},\quad
 \norm{\Delta_{\Sp}\psi_{j,h}}_{L^2(\Sp)}^2\asymp h^{d-5},
 \quad \norm{D^k\psi_{j,h}}_\infty\le C_kh^{-k},
\end{equation}
the derivatives being taken in the chart. All but the lower bound in
the middle relation are scaling; for that
one, integration by parts on the closed manifold $\Sp$ gives
$\norm{\Delta\psi}_2\norm\psi_2\ge\norm{\nabla\psi}_2^2$, and
$\norm{\nabla\psi_{j,h}}_2^2\asymp h^{d-3}$ by scaling.

Choose a nonzero nonnegative radial profile
$\eta\in C_c^\infty((1/4,1/2))$. For $\theta\in\{0,1\}^{N_h}$ and a
small fixed amplitude $\kappa>0$, to be chosen in Section
\ref{sec:membership}, prescribe the solutions
\begin{equation}\label{eq:solutions}
 u_\theta=u_*+v_\theta,\qquad
 v_\theta(\rho\omega)=
 \kappa h^{\alpha+2}\eta(\rho)
             \sum_{j=1}^{N_h}\theta_j\psi_{j,h}(\omega).
\end{equation}
The perturbation $v_\theta$ vanishes for $\rho<1/4$ and for
$\rho>1/2$, so it is smooth on $\R^d$ and $u_\theta=0$ on $\partial D$.
The exponent $\alpha+2$ is dictated by the parameter class: two angular
derivatives of $v_\theta$ will become $f_\theta-1$, which must have
$C^\alpha$ norm of order $\kappa$. \review{Clarify that the last statement only serves intuition and 'become' does not mean 'equal'.} Since the bumps have disjoint
supports, for every fixed integer $k\ge0$,
\begin{equation}\label{eq:vderivatives}
 \norm{D^kv_\theta}_\infty\le C_k\kappa h^{\alpha+2-k},
 \qquad \norm{v_\theta}_2^2\le C\kappa^2h^{2\alpha+4},
\end{equation}
the second because $\sum_j\norm{\psi_{j,h}}_2^2\asymp N_hh^{d-1}\asymp1$.
In particular $\norm{D^2v_\theta}_\infty\le C\kappa h^\alpha\le C\kappa$,
so for small $\kappa$ the Hessians satisfy \eqref{eq:geometry} on
$\overline D$, and $\nabla u_\theta\cdot\nu=\nabla u_*\cdot\nu=1/d>0$
on $\partial D$. Lemma \ref{lem:transport} applies on $\Omega=D$, and
we define $f_\theta=T_{u_\theta}^{-1}1=1+q_\theta$. Since
$T_{u_\theta}1=\Delta u_\theta=1+\Delta v_\theta$, the deviation
$q_\theta$ solves the transport equation
\begin{equation}\label{eq:qequation}
 T_{u_\theta}q_\theta=-\Delta v_\theta .
\end{equation}

\subsection{Membership in the parameter class}\label{sec:membership}

We claim that, for every fixed integer $k\ge0$ (only
$k\le\lfloor\alpha\rfloor+1$ is needed),
\begin{equation}\label{eq:qderivatives}
 \norm{D^kq_\theta}_\infty\le C_k\kappa h^{\alpha-k}.
\end{equation}
For $k=0$, \eqref{eq:qequation}, the supremum bound
\eqref{eq:inversebounds} and \eqref{eq:vderivatives} give
$\norm{q_\theta}_\infty\le c_0^{-1}\norm{\Delta v_\theta}_\infty
\le C\kappa h^\alpha$. For the induction step, differentiate
\eqref{eq:qequation} $k$ times. The result is a transport equation for
the symmetric tensor $D^kq_\theta$ along the same flow, in which the
first-order part is $\nabla u_\theta\cdot\nabla(D^kq_\theta)$ and the
zero-order principal coefficient is $\Delta u_\theta$ plus
$D^2u_\theta$ acting in each of the $k$ tensor slots (one such term
for each way of placing exactly one derivative on $\nabla u_\theta$);
this symmetric coefficient is bounded below by $c_0+km$. The main
forcing is $D^k\Delta v_\theta$, of size $C_k\kappa h^{\alpha-k}$. The
remaining terms come from placing at least two derivatives on
$\nabla u_\theta$, or at least one on $\Delta u_\theta$; since $u_*$ is
quadratic, only $v_\theta$ contributes, and they are products of the
forms
\[
 D^j\Delta v_\theta\,D^{k-j}q_\theta\quad(1\le j\le k),
 \qquad
 D^{j+1}v_\theta\,D^{k+1-j}q_\theta\quad(2\le j\le k).
\]
By \eqref{eq:vderivatives} and the induction hypothesis (which applies
since $k-j\le k-1$ and $k+1-j\le k-1$), each is bounded by
$C_k\kappa^2h^{2\alpha-k}\le C_k\kappa h^{\alpha-k}$. The tensor
version of the flow bound, proved exactly as
\eqref{eq:gradientinverse} with the Frobenius norm of the tensor in
place of $|w|$, gives \eqref{eq:qderivatives}. \review{The entire argument is a little  fuzzy. Please write down the transport equation for the tensors, at least symbolically, and argue then with precise attributions to the terms.}

From \eqref{eq:qderivatives} to the H\"older--Zygmund norm: if
$\alpha$ is not an integer, the supremum bound for $D^{\lfloor\alpha\rfloor}
q_\theta$ and the Lipschitz bound coming from
$D^{\lfloor\alpha\rfloor+1}q_\theta$ interpolate to the H\"older
seminorm of exponent $\alpha-\lfloor\alpha\rfloor$, with the powers of
$h$ cancelling exactly; if $\alpha$ is an integer, bounded classical
derivatives of order $\alpha$ control the Zygmund norm. (These
classical bounds on $\overline D$ pass to the quotient norm of Section
1.1 through a bounded extension to $\R^d$, e.g. the reflection
extension \cite{Hestenes,Seeley} across the smooth boundary.) Thus
\begin{equation}\label{eq:qholder}
 \norm{f_\theta-1}_{C^\alpha(D)}\le C\kappa .
\end{equation}
Fix $\kappa>0$ so small that the right-hand side is at most $r$ and
the Hessian condition above holds. Then all $f_\theta$ belong to
$\mathcal F(\alpha,r)$, in particular they are positive. By
construction $\diver(f_\theta\nabla u_\theta)=1$ in $D$ and
$u_\theta|_{\partial D}=0$, so uniqueness of the weak solution gives
$u_\theta=u_{f_\theta}$ exactly.

\subsection{Separation of the alternatives}

Let $T_*=T_{u_*}=d^{-1}x\cdot\nabla+1$ \review{Explain that $1$ denotes the multiplication operator with $1$, so the identity} be the transport operator of
the reference solution, and let $w_\theta=-T_*^{-1}\Delta v_\theta$ be
the linearization of $q_\theta$. Since $T_{u_\theta}=T_*+T_{v_\theta}$,
\eqref{eq:qequation} gives the exact identity
$T_*q_\theta=-\Delta v_\theta-\diver(q_\theta\nabla v_\theta)$,
so $q_\theta-w_\theta=-T_*^{-1}\diver(q_\theta\nabla v_\theta)$, and
\eqref{eq:inversebounds}, \eqref{eq:vderivatives} and
\eqref{eq:qderivatives} imply
\begin{equation}\label{eq:lowerremainder}
 \norm{q_\theta-w_\theta}_2
 \le C\norm{\diver(q_\theta\nabla v_\theta)}_\infty
 \le C\big(\norm{\nabla q_\theta}_\infty\norm{\nabla v_\theta}_\infty
 +\norm{q_\theta}_\infty\norm{\Delta v_\theta}_\infty\big)
 \le C\kappa^2h^{2\alpha}.
\end{equation}
The inverse of $T_*$ acts along rays: solving
$d^{-1}\rho\,\partial_\rho w+w=z$, i.e. $\partial_\rho(\rho^dw)=d\rho^{d-1}z$,
gives $T_*^{-1}z(\rho\omega)=d\rho^{-d}\int_0^\rho s^{d-1}z(s\omega)\,ds$.
In polar coordinates
$\Delta(\eta\psi)=(\eta''+(d-1)\rho^{-1}\eta')\psi+\rho^{-2}\eta\,\Delta_{\Sp}\psi$,
and the radial part integrates to
$\int_0^\rho s^{d-1}(\eta''+(d-1)s^{-1}\eta')\,ds=\rho^{d-1}\eta'(\rho)$,
which vanishes for $\rho>1/2$, outside the support of $\eta$.
Therefore, for $1/2<\rho<1$,
\begin{equation}\label{eq:tail}
 w_\theta(\rho\omega)
 =-d\kappa h^{\alpha+2}\rho^{-d}A_\eta
       \Delta_{\Sp}\sum_j\theta_j\psi_{j,h}(\omega),
 \qquad A_\eta=\int_{1/4}^{1/2}s^{d-3}\eta(s)\,ds>0:
\end{equation}
the two angular derivatives convert the amplitude $h^{\alpha+2}$ into
$h^\alpha$. If the Hamming distance between $\theta,\theta'$ is
at least $cN_h$, the disjointness of the supports of the
$\Delta_{\Sp}\psi_{j,h}$, \eqref{eq:angularbumps} and
$N_hh^{d-1}\asymp1$ give
$\norm{w_\theta-w_{\theta'}}_2^2
 \ge c\kappa^2 h^{2\alpha+4}\cdot cN_h\,h^{d-5}
 \ge c'\kappa^2h^{2\alpha}$. By \eqref{eq:lowerremainder} and the
triangle inequality, for all sufficiently small $h$ the corresponding
conductivities satisfy
\begin{equation}\label{eq:separation}
 \norm{f_\theta-f_{\theta'}}_{L^2(D)}=\norm{q_\theta-q_{\theta'}}_2
 \ge c''\kappa h^\alpha .
\end{equation}
\review{The argument would be easier to follow when stating a lower bound for $\|w_\theta-w_{\theta'}\|_2$, not its square so that the quadratically small term in (6.7) is immediately seen to be of smaller order. Write explicitly $c''>0$ and add an intermediate step.}

\subsection{Information bound and Fano's inequality}

The Varshamov--Gilbert bound \cite[Lemma 2.9]{Tsybakov} provides a
subset $\Theta_h\subset\{0,1\}^{N_h}$ with pairwise Hamming distances
at least $N_h/8$ and $\log|\Theta_h|\ge cN_h\ge ch^{-(d-1)}$. For the
Gaussian white noise model \eqref{eq:model}, the Kullback--Leibler
divergence between two shifts is the exact Gaussian expression
\cite[equation (6.16)]{GN}, so, writing $P_{\mathbf 1}$ for the law
of $Y_n$ under $f=\mathbf 1$, \eqref{eq:vderivatives} gives
\begin{equation}\label{eq:KL}
 \KL(P_{f_\theta},P_{\mathbf 1})
       =\frac n2\norm{u_\theta-u_*}_2^2
       =\frac n2\norm{v_\theta}_2^2
       \le Cn\kappa^2h^{2\alpha+4}.
\end{equation}
Choose $h=c_*n^{-1/(2\alpha+d+3)}$ with $c_*>0$ small. Then
$nh^{2\alpha+4}=c_*^{2\alpha+d+3}h^{-(d-1)}$, so \eqref{eq:KL} is at
most $\tfrac1{16}\log|\Theta_h|$ for small $c_*$.

Let $\vartheta$ be uniformly distributed on $\Theta_h$ and $\bar P$
the mixture of the $P_{f_\theta}$. For any probability measure $Q$
with respect to which the $P_{f_\theta}$ have finite Kullback--Leibler
divergence, writing $dP_{f_\theta}/d\bar P$ as
$(dP_{f_\theta}/dQ)/(d\bar P/dQ)$ inside the logarithm and averaging
over $\theta$ shows that the mutual information between $\vartheta$
and $Y_n$ satisfies
\begin{align*}
 I(\vartheta;Y_n)
 &=|\Theta_h|^{-1}\textstyle\sum_\theta\KL(P_{f_\theta},\bar P)
 =|\Theta_h|^{-1}\sum_\theta\KL(P_{f_\theta},Q)-\KL(\bar P,Q)\\
 &\le|\Theta_h|^{-1}\textstyle\sum_\theta\KL(P_{f_\theta},Q)
 \le\tfrac1{16}\log|\Theta_h|\qquad\text{for }Q=P_{\mathbf 1}.
\end{align*}
Fano's inequality \cite[Lemma 2.10 and Corollary 2.6]{Tsybakov} therefore
bounds the average misclassification probability of every classifier
$\widehat\vartheta=\widehat\vartheta(Y_n)$ with values in $\Theta_h$
below by
\[
 \frac1{|\Theta_h|}\sum_{\theta\in\Theta_h}
 P_{f_\theta}(\widehat\vartheta\ne\theta)\ \ge\
 1-\frac{\frac1{16}\log|\Theta_h|+\log 2}{\log|\Theta_h|}
 \ge \frac12
\]
for large $n$. Given an arbitrary estimator $\widehat f_n$, let
$\widehat\vartheta$ be the index of a nearest $f_\theta$ in $L^2(D)$.
By \eqref{eq:separation}, $\widehat\vartheta\ne\theta$ forces
$\norm{\widehat f_n-f_\theta}_2\ge c''\kappa h^\alpha/2$, so
\[
 \max_{\theta\in\Theta_h}\E_{f_\theta}\norm{\widehat f_n-f_\theta}_2^2
 \ge\frac{(c''\kappa h^\alpha)^2}4\cdot
 \frac1{|\Theta_h|}\sum_\theta P_{f_\theta}(\widehat\vartheta\ne\theta)
 \ge c\kappa^2h^{2\alpha}=c\,n^{-\frac{2\alpha}{2\alpha+d+3}} .
\]
Since all $f_\theta$ lie in $\mathcal F(\alpha,r)$, this is the claimed
lower bound and completes the proof of Theorem \ref{thm:main}.

\section{Remarks on the mechanism}\label{sec:remarks}

\begin{remark}
The two halves of the proof explain the denominator
$2\alpha+d+3=2\alpha+(d-1)+4$ from opposite directions: the alternatives oscillate in the $d-1$
angular directions and their solutions are smaller than their
conductivities by two angular derivatives, while integration along the
gradient flow averages the noise in the one flow direction but not in
the $d-1$ transverse ones. No boundary values of $f$ were assumed, and
the linearized alternatives \eqref{eq:tail} are nonconstant up to the
boundary. Besides the classical embedding, $\alpha>1$ enters through
$\sum_i2^{(1-\alpha)i}<\infty$ and $2\alpha>\alpha+1$ in the proof of
Lemma \ref{lem:comm}, and through the squared nonlinear remainder
\eqref{eq:remainder} being $o(h^{2\alpha})$.
\end{remark}

\begin{thebibliography}{99}

\bibitem{GiNi}
M. Giordano and R. Nickl,
Consistency of Bayesian inference with Gaussian process priors in an
elliptic inverse problem,
Inverse Problems 36 (2020), 085001.

\bibitem{GT}
D. Gilbarg and N. S. Trudinger,
Elliptic Partial Differential Equations of Second Order,
Springer, Berlin, 2001.

\bibitem{GN}
E. Gin\'e and R. Nickl,
Mathematical Foundations of Infinite-Dimensional Statistical Models,
Cambridge University Press, Cambridge, 2016.

\bibitem{Hestenes}
M. R. Hestenes,
Extension of the range of a differentiable function,
Duke Math. J. 8 (1941), 183--192.

\bibitem{KSV}
G. Koers, B. Szab\'o and A. van der Vaart,
Linear methods for non-linear inverse problems,
arXiv:2411.19797, 2024.

\bibitem{Nickl23}
R. Nickl,
Bayesian Non-linear Statistical Inverse Problems,
EMS Press, Berlin, 2023.

\bibitem{NP}
R. Nickl and G. P. Paternain,
On some information-theoretic aspects of non-linear statistical
inverse problems,
Proc. Int. Congress of Mathematicians 2022, Vol. 7, EMS Press, 2023,
5516--5538.

\bibitem{NvdGW}
R. Nickl, S. van de Geer and S. Wang,
Convergence rates for penalized least squares estimators in PDE
constrained regression problems,
SIAM/ASA J. Uncertain. Quantif. 8 (2020), 374--413.

\bibitem{Richter}
G. R. Richter,
An inverse problem for the steady state diffusion equation,
SIAM J. Appl. Math. 41 (1981), 210--221.

\bibitem{Seeley}
R. T. Seeley,
Extension of $C^\infty$ functions defined in a half space,
Proc. Amer. Math. Soc. 15 (1964), 625--626.

\bibitem{Triebel}
H. Triebel,
Theory of Function Spaces II,
Birkh\"auser, Basel, 1992.

\bibitem{Tsybakov}
A. B. Tsybakov,
Introduction to Nonparametric Estimation,
Springer, New York, 2009.

\end{thebibliography}

\end{document}
